\subsection{Полное внутреннее отражение}

%%%%%%%%%%%%%%%%%%%%%%%%%%%%%%%%%%%%%%%%%%%%%%%%%%%%
%%%%%%%%%%%%%%%%%%%%%%%%%%%%%%%%%%%%%%%%%%%%%%%%%%%%
%%%%%%%%%%%%%%%%%%%%%%%%%%%%%%%%%%%%%%%%%%%%%%%%%%%%
 \z[{\begin{wrapfigure}[5]{r}{0.25\textwidth}%
		\vspace{-10pt}%
		\begin{flushright}%
 			\includegraphics[width=0.25\textwidth]{4.4.1}%
		\end{flushright}%
	\end{wrapfigure}%
	}]%
{На призму, сечение которой имеет вид равнобедренного прямоугольного треугольника, перпендикулярно нижней грани падает луч от лазерной указки. Каким должен быть показатель преломления $n$ материала, из которого сделана призма, чтобы свет от указки вышел из призмы наружу только через эту же грань?}
%
{$n > \sqrt{2} $}

%%%%%%%%%%%%%%%%%%%%%%%%%%%%%%%%%%%%%%%%%%%%%%%%%%%%
 \z{Снаружи круглого прозрачного стержня вблизи от центра его торца помещён точечный источник света. При каких значениях показателя преломления материала стержня $n$ свет не будет выходить через его боковую поверхность?}
%
{$n > \sqrt{2} $}

%%%%%%%%%%%%%%%%%%%%%%%%%%%%%%%%%%%%%%%%%%%%%%%%%%%%
 \z{На поверхности водоёма, имеющего глубину $H = 3.3\,\text{м}$, плавает фанерный круг радиусом $r = 3\,\text{м}$. На оси круга расположен точечный источник света, высота которого над поверхностью круга может меняться. Чему равен максимальный радиус тени круга на дне $R$, если показатель преломления воды $n = 1.33$?}
%
{$R = r + \frac{H}{\sqrt{n^{2} - 1}} = 6.8\,\text{м}$}

%%%%%%%%%%%%%%%%%%%%%%%%%%%%%%%%%%%%%%%%%%%%%%%%%%%%
 \z{Снаружи круглого прозрачного стержня вблизи от центра его торца помещён точечный источник света. При каких значениях показателя преломления материала стержня $n$ свет не будет выходить через его боковую поверхность?}
%
{$n > \sqrt{2}$}

%%%%%%%%%%%%%%%%%%%%%%%%%%%%%%%%%%%%%%%%%%%%%%%%%%%%
 \z{На горизонтальном дне водоёма лежит монета радиуса $r = 2\,\text{см}$. На каком максимальном расстоянии от монеты надо поместить в воде плоский экран радиуса $R = 5\,\text{см}$, чтобы монету нельзя было обнаружить из воздуха при спокойной поверхности воды? Показатель преломления воды $n = \frac{4}{3}$.}
%
{$L = (R - r) \sqrt{n^{2} - 1} = 2.6\,\text{см}$}

%%%%%%%%%%%%%%%%%%%%%%%%%%%%%%%%%%%%%%%%%%%%%%%%%%%%
 \z{В стекле с показателем преломления $n_1 = 1.5$ имеется сферическая полость радиуса $R = 4.5\,\text{см}$, заполненная водой. Показатель преломления воды $n_2 = \frac{4}{3}$. На полость падает широкий пучок параллельных световых лучей. Определить радиус $r$ пучка световых лучей, которые проникают в полость.}
%
{$r = R \cdot \frac{n_{2}}{n_{1}}$}

%%%%%%%%%%%%%%%%%%%%%%%%%%%%%%%%%%%%%%%%%%%%%%%%%%%%
 \z[{\begin{wrapfigure}[5]{r}{0.25\textwidth}%
		\vspace{-10pt}%
		\begin{flushright}%
 			\includegraphics[width=0.25\textwidth]{4.4.7}%
		\end{flushright}%
	\end{wrapfigure}%
	}]%
{Луч света, лежащий в плоскости рисунка, падает на боковую грань $AB$ призмы, имеющей при вершине угол $90\degree$. В каких пределах лежат возможные значения угла падения $\alpha$, если известно, что луч выходит из боковой грани $AC$? Показатель преломления призмы $n = 1.25$.}
%
{$\sin{\alpha} > \sqrt{n^{2} - 1}$}



